\documentclass[conference]{IEEEtran}

\IEEEoverridecommandlockouts

\usepackage{cite}

\usepackage{amsmath,amssymb,amsfonts}

\usepackage{algorithmic}

\usepackage{graphicx}

\usepackage{textcomp}

\usepackage{xcolor}

\usepackage{hyperref}

\usepackage[ruled,vlined]{algorithm2e}


\def\BibTeX{{\rm B\kern-.05em{\sc i\kern-.025em b}\kern-.08em
    T\kern-.1667em\lower.7ex\hbox{E}\kern-.125emX}}

\setlength{\parindent}{0pt}
\begin{document}


\title{Performance Analysis of Sparse Matrix-Vector Multiplication (SpMV)}


\author{\IEEEauthorblockN{Michael Bernasconi}
\textit{Student ID: 267681}\\
\textit{michael.bernasconi@studenti.unitn.it} \\
\url{https://github.com/Michael-Bernasconi/Sparse-Matrix-Vector-Multiplication}
\IEEEauthorblockA{\textit{Department of Information Engineering and Computer Science} \\
\textit{University of Trento}\\
}
}


\maketitle


\begin{abstract}

Sparse Matrix-Vector multiplication (SpMV) is a fundamental cornerstone of numerical computing; its efficiency on GPUs is historically limited by irregular memory accesses and load imbalance. This study analyzes the performance of custom implementations (COO, CSR, CSR-Vector) by comparing them with the cuSPARSE library on NVIDIA Ampere architecture (A30). A novel evaluation is proposed that correlates matrix topology with hardware efficiency: the results demonstrate that the CSR-Vector approach, through warp-level parallel reduction, mitigates computational divergence of irregular matrices, achieving a peak of 95.73 GFLOPS and saturating up to 390.83 GB/s of effective bandwidth. A distinctive contribution of this work is the analysis of Time To Solution (TTS), which reveals that I/O parsing costs and PCIe transfers dominate the program lifecycle.

\end{abstract}


\begin{IEEEkeywords}

SpMV, CUDA, Sparse Matrix Vector Multiplication, Load Balancing, NVIDIA Ampere, Memory Bandwidth, Time To Solution (TTS), Performance Evaluation.

\end{IEEEkeywords}


\section{Introduction}

The Sparse Matrix-Vector Multiplication (SpMV) operation, defined as $y = Ax$, constitutes a critical computational kernel for High Performance Computing (HPC) \cite{bGao2024}.

Due to its low arithmetic intensity, quantifiable as $2 \times nnz$ total operations, SpMV is an inherently memory-bound operation.

Its efficiency is further hindered by the irregular nature of memory accesses, which strongly depend on the sparsity structure of the matrix.

This study analyzes the impact of storage formats and parallelization strategies on NVIDIA Ampere architecture, comparing custom CUDA implementations (COO, CSR, and CSR-Vector) with the NVIDIA cuSPARSE library.\\

The investigation evaluates performance through key metrics: kernel execution time, throughput (GFLOPS), effective bandwidth (BW), and overall \textit{Time To Solution} (TTS).

This study correlates structural characteristics of the matrix with hardware performance, offering an experimental guide for selecting the optimal format and kernel in order to mitigate computational and system bottlenecks.


\section{Methodology}

This section describes the test environment, the state-of-the-art context for SpMV optimizations, the matrix formats examined, and the rigorous methodology adopted for performance measurement and validation.


\subsection{Hardware and Software Environment}

Computations were performed on the UniTN cluster (node edu01) on the edu-short partition. Hardware specifications are summarized in Table \ref{tab:hardware}

\begin{table}[htbp]

\caption{Hardware Environment Specifications}

\label{tab:hardware}

\begin{center}

\footnotesize

\setlength{\tabcolsep}{3pt}

\begin{tabular}{|l|p{3.5cm}|p{3.5cm}|}

\hline

Feature & Host (CPU) & Device (GPU) \\

\hline

Model & Intel Xeon Silver 4309Y & NVIDIA A30 \\

Architecture & Ice Lake (x86\_64) & Ampere (Compute Cap. 8.0) \\

Core / SM & 16 Core / 32 Thread (2 Sock.) & 56 SM (3584 CUDA Core) \\

Clock Frequency & 2.80 GHz (Base) / 3.60 GHz & 1.44 GHz (Boost) \\

FP32 Performance & 1.84 TFLOPS & 10.3 TFLOPS \\

Memory Bandwidth & 102.4 GB/s & 933.1 GB/s (HBM2) \\

Total Global Memory & N/A & 24 GB \\

Shared Memory & N/A & 48-100 KiB / SM \\

Warp Size & N/A & 32 \\

\hline

\end{tabular}

\end{center}

\end{table}


\subsection{State of the Art in SpMV Optimization}

SpMV is a critical bottleneck in HPC due to bandwidth saturation and irregular memory access patterns. Although the CSR format was originally designed to minimize memory footprint in sequential contexts, GPU architectures require strategies to mitigate warp divergence and improve memory coalescing \cite{bChu2023}. Modern approaches, such as cuSPARSE and warp-based kernels, use cooperative strategies to handle row variance \cite{bBell2009}.


\subsection{CPU Baseline and GPU Implementations}

CPU Baseline: The reference computation on the CPU was handled through a parallel implementation based on OpenMP. This computational phase is not intended as a performance benchmark, but is used for validation and correctness checking of results subsequently computed in parallel on the GPU.


In GPU computation, SpMV performance critically depends on the interaction between the sparse matrix memory layout and the adopted execution model. To analyze this dual aspect (data structure and computational logic), our investigation focuses on the following approaches:


\subsubsection{COO Format and Base Kernel}

The Coordinate (COO) format explicitly stores the sparse matrix through three arrays of size equal to the number of non-zero elements (NNZ): row indices, column indices, and values. As described in Algorithm \ref{alg:coo}, the kernel maps one thread per element.


\begin{algorithm}[ht]
    \footnotesize
\caption{SpMV Kernel: COO Base (Atomic-based)}\label{alg:coo_new}
\begin{algorithmic}[1]
\REQUIRE $val, row\_ind, col\_ind$ (arrays of size $nnz$), $x$ (dense vector)
\ENSURE $y$ (output vector)
\STATE $i \gets$ Global thread index
\IF{$i < nnz$}
\STATE $v \gets val[i]$
\STATE $r \gets row\_ind[i]$
\STATE $c \gets col\_ind[i]$
\STATE $\text{AtomicAdd}(y[r], v \times x[c])$ \COMMENT{Handle race conditions on $y$}
\ENDIF
\end{algorithmic}
\end{algorithm}


The Base COO kernel adopts a direct mapping strategy: it launches a single thread for each non-zero element. This approach guarantees excellent load balancing since it is insensitive to row distribution. However, since multiple threads may update the same row concurrently, it requires \texttt{atomicAdd} to resolve race conditions, limiting throughput.


\subsubsection{CSR Format and Base Kernel (Scalar)}

The Compressed Sparse Row (CSR) format optimizes memory footprint by compressing row indices into a \texttt{row\_ptr} vector of size $M+1$, following the logic of Algorithm \ref{alg:csr}.

\begin{algorithm}[h]
    \footnotesize
\caption{SpMV Kernel: CSR Scalar}\label{alg:csr_new}
\begin{algorithmic}[1]
\REQUIRE $row\_ptr$ (size $M+1$), $col\_idx, val$ (size $nnz$), $x$ (dense vector)
\STATE $r \gets$ Assigned row index per thread
\IF{$r < M$}
\STATE $sum \gets 0$
\STATE $row\_start \gets row\_ptr[r]$
\STATE $row\_end \gets row\_ptr[r+1]$
\FOR{$j$ from $row\_start$ to $row\_end - 1$}
\STATE $sum \gets sum + val[j] \times x[col\_idx[j]]$
\ENDFOR
\STATE $y[r] \gets sum$
\ENDIF
\end{algorithmic}
\end{algorithm}


The Base CSR kernel maps one thread per row. This eliminates atomic operations but suffers from \textit{warp divergence} and non-coalesced memory accesses when rows have significantly different lengths.


\subsubsection{Optimized CSR-Vector Kernel}

To overcome inefficiencies of scalar CSR, the CSR-Vector kernel (Algorithm \ref{alg:csr_vector}) redefines execution granularity by assigning an entire warp (32 threads) to process a single row. This enables coalesced memory access and uses shared memory for parallel reduction.

\begin{algorithm}[h]
    \footnotesize
\caption{SpMV Kernel: CSR-Vector (Warp-level Parallelism)}\label{alg:csr_vector_new}
\begin{algorithmic}[1]
\REQUIRE $row\_ptr, col\_idx, val, x$
\STATE $lane \gets$ Thread index within the Warp ($0 \dots 31$)
\STATE $row \gets$ Assigned row index per Warp
\IF{$row < M$}
\STATE $sum \gets 0$
\STATE $start \gets row\_ptr[row]$
\STATE $end \gets row\_ptr[row+1]$
\FOR{$j = start + lane$ \TO $end-1$ \textbf{step} 32}
\STATE $sum \gets sum + val[j] \times x[col\_idx[j]]$ \COMMENT{Coalesced access}
\ENDFOR
\STATE $shared\_data[threadID] \gets sum$
\STATE $\text{WarpSync}()$
\STATE \COMMENT{Parallel Tree Reduction}
\FOR{$offset \in \{16, 8, 4, 2, 1\}$}
\IF{$lane < offset$}
\STATE $shared\_data[tid] \gets shared\_data[tid] + shared\_data[tid + offset]$
\ENDIF
\STATE $\text{WarpSync}()$
\ENDFOR
\IF{$lane = 0$}
\STATE $y[row] \gets shared\_data[threadID]$
\ENDIF
\ENDIF
\end{algorithmic}
\end{algorithm}

Finally, all custom implementations were compared with the cuSPARSE library, which serves as a production-level baseline for NVIDIA Ampere architecture.

Benchmark accuracy is ensured by the following methodological choices:


\begin{itemize}

    \item \textit{Output Consistency:} To prevent accumulation of incorrect values across iterations (a critical phenomenon in COO format using \texttt{atomicAdd}), the destination vector is explicitly reset before each kernel launch.

    

    \item \textit{Variability Mitigation and Measurement:} Pure computation time (\textit{Kernel Time}) is evaluated, ignoring initial I/O (parsing) and memory allocation costs. To obtain stable measurements, each test is divided into two phases:

    \begin{itemize}

        \item Warm-up (20 iterations): Necessary to stabilize hardware and bring GPU frequencies to steady state.

        \item Benchmark (500 iterations): Actual performance measurement.

    \end{itemize}

    The entire procedure is repeated for 5 independent \textit{runs}. The average time recorded during benchmark phases (accompanied by standard deviation to prove stability) becomes the baseline parameter on which all final performance metrics are computed: Kernel Time, Time to Solution (TTS), Throughput (GFLOPS), and Effective Bandwidth (BW).

\end{itemize}


\subsection{Performance Metrics}Implementations are evaluated using the following metrics and equations:\begin{itemize} \itemsep0em\item Kernel Time ($T_{avg}$): Recorded using \texttt{cudaEvents} to sample exclusively the kernel time window on VRAM, avoiding CPU-related latencies.\item Time To Solution (TTS): Measured using native C libraries (\texttt{my\_time\_lib.h}), capturing the full program lifecycle from startup to validation completion. Includes Matrix Market file parsing costs, \texttt{cudaMalloc}, and \texttt{cudaMemcpy} overhead.


\item Computational Throughput (GFLOPS): The standard assumes 2 float operations (FMA) per non-zero element: $GFLOPS = \frac{2 \times \text{NNZ}}{T_{avg} \times 10^9}$.


\item Effective Bandwidth (BW) (GB/s): Being a \textit{memory bound} operation, bandwidth varies with format.

For CSR: $BW_{CSR} = \frac{4 (M + 1 + \text{nnz}) + 4 (\text{nnz} + N + M)}{T_{avg} \times 10^9}$.

For COO: $BW_{COO} = \frac{4 (2 \times \text{nnz}) + 4 (\text{nnz} + N + M)}{T_{avg} \times 10^9}$, where factor 4 represents bytes for \texttt{int32} and \texttt{float32}.


\item Cache Statistics (CPU): Profiling of the hierarchical cache system was conducted using Valgrind: $\text{Miss D1 (\%)} = \left( \frac{\text{Miss D1}}{\text{Total D1 Accesses}} \right) \times 100$. It allows quantifying access efficiency and theoretical bandwidth collapse.

\end{itemize}


\subsection{Correctness and Validation}

Validation is performed by comparing GPU results with the CPU baseline using a relative/absolute numerical tolerance ($10^{-3}$), necessary to handle non-associativity of parallel \texttt{Float32} arithmetic. In a concurrent execution context, non-deterministic operation ordering introduces unavoidable rounding errors compared to sequential execution. Comparison is performed via the \texttt{validate\_results()} method: if divergence between GPU and CPU output exceeds the threshold, the experiment reports a FAIL state, ensuring mathematical integrity of results.


\section{Dataset}

The experimental evaluation is conducted on a selection of 10 sparse matrices from the \textit{\href{https://sparse.tamu.edu/}{SuiteSparse Matrix Collection}}, whose details are reported in Table \ref{tab:datasets}.
These datasets were chosen to represent a wide variety of application domains and structural characteristics, ranging from perfectly balanced circuit simulations to extremely irregular power-law graphs. Notably, this selection aligns with recent high-performance computing literature, as these matrices are extensively utilized for benchmarking in the work of Chu et al. \cite{bChu2023}.

\begin{table}[htbp]
\caption{Structural characteristics of selected SuiteSparse matrices}
\label{tab:datasets}
\begin{center}
\footnotesize
\setlength{\tabcolsep}{4pt}
\begin{tabular}{|l|r|r|r|l|}
\hline
\textbf{Matrix Name} & \textbf{NNZ} & \textbf{Avg NNZ/R} & \textbf{Std. Dev.} & \textbf{Structure} \\ \hline
ASIC\_680ks     & 1,693,767  & 2.48  & 4.16    & Regular \\ \hline
bone010         & 47,851,783 & 48.50 & 7.62    & Dense, Symmetric \\ \hline
boyd2           & 1,500,397  & 3.22  & 194.91  & Irregular \\ \hline
FullChip        & 26,621,983 & 8.91  & 1,806.80& Extreme Imbalance \\ \hline
Ga41As41H72     & 18,488,476 & 68.96 & 105.39  & Dense, Symmetric \\ \hline
ldoor           & 42,493,817 & 44.63 & 14.77   & Structured \\ \hline
rajat31         & 20,316,253 & 4.33  & 1.11    & Very Regular \\ \hline
Rucci1          & 7,791,168  & 3.94  & 0.34    & Rectangular \\ \hline
Si41Ge41H72     & 15,011,265 & 80.86 & 126.97  & Dense \\ \hline
webbase-1M      & 3,105,536  & 3.11  & 25.35   & Power-law \\ \hline
\end{tabular}
\end{center}
\end{table}
\section{Results}

This section presents experimental data obtained from evaluation of SpMV implementations. 


\subsection{Computational Throughput (GFLOPS)}

Maximum throughput (Figure \ref{fig:gflops}) was recorded by the \texttt{CSR-Vector} kernel on matrix \texttt{bone010} with 95.73 GFLOPS. The \texttt{Base CSR} implementation showed severe inefficiencies on irregular matrix \texttt{FullChip}, dropping to 0.1989 GFLOPS, while \texttt{cuSPARSE} maintained performance of 15.68 GFLOPS.


\begin{figure}[htbp]

\centerline{\includegraphics[width=\linewidth]{C:/Users/micha/Documents/GitHub/Sparse-Matrix-Vector-Multiplication/results/plots/1_del_gflops.pdf}}

\caption{Computational Throughput (GFLOPS) for tested implementations.}

\label{fig:gflops}

\end{figure}


\subsection{Kernel Execution Time}

Average times (Figure \ref{fig:kernel_time}) confirm throughput dynamics. For matrix \texttt{FullChip}, execution decreases from 269 ms (\texttt{Base CSR}) to 14.1 ms (\texttt{CSR-Vector}). The \texttt{cuSPARSE} library proves fastest in complex scenarios, recording 3.4 ms on the same matrix.


\begin{figure}[htbp]

\centerline{\includegraphics[width=\linewidth]{C:/Users/micha/Documents/GitHub/Sparse-Matrix-Vector-Multiplication/results/plots/1_del_kernel_time.pdf}}

\caption{Average kernel time in seconds with standard deviation.}

\label{fig:kernel_time}

\end{figure}


\subsection{Time To Solution (TTS)}

Figure \ref{fig:tts} illustrates Total Time To Solution (TTS) of the application. The logarithmic scale graph highlights how inclusion of all preparatory phases tends to level performance differences previously measured on individual kernels. On matrix \texttt{ASIC\_680ks}, for example, CPU implementation takes about 40.38 s, while the four GPU variants (COO, CSR, CSR-Vector, and cuSPARSE) all converge almost identically between 8.41 s and 8.59 s.


\begin{figure}[htbp]

\centerline{\includegraphics[width=\linewidth]{C:/Users/micha/Documents/GitHub/Sparse-Matrix-Vector-Multiplication/results/plots/1_del_tts.pdf}}

\caption{Total Time To Solution (TTS) in seconds measured on logarithmic scale for different implementations.}

\label{fig:tts}

\end{figure}


\subsection{Effective Memory Bandwidth}

Peak measured value (Figure \ref{fig:bandwidth}) is 390.83 GB/s (\texttt{CSR-Vector} on \texttt{bone010}), highlighting a practical limit compared to theoretical bandwidth of NVIDIA A30 (933.1 GB/s). CPU implementations reach a maximum of 13.99 GB/s (\texttt{boyd2}).


\begin{figure}[htbp]

\centerline{\includegraphics[width=\linewidth]{C:/Users/micha/Documents/GitHub/Sparse-Matrix-Vector-Multiplication/results/plots/1_del_bandwidth.pdf}}

\caption{Effective Memory Bandwidth (GB/s). Theoretical peak of A30 is 933.1 GB/s.}

\label{fig:bandwidth}

\end{figure}


\subsection{Cache Profiling}

Analysis via Valgrind detected an extremely low D1 Miss Rate ($<0.3\%$) for profilable datasets (\textit{ASIC\_680ks, boyd2, Rucci1, webbase-1M}), indicating good local data reuse on vectors. Conversely, for large matrices (\textit{bone010, FullChip, ldoor, rajat31}), simulation exceeded the 5-minute limit imposed by the cluster scheduler. This timeout is attributable to high computational overhead of Valgrind when tracing indirect memory accesses typical of SpMV. On datasets with millions of non-zero elements, simulation latency of sparse accesses saturates available time before kernel completion, preventing exhaustive measurement of memory hierarchy.


\subsection{Numerical Validation}

The validation procedure systematically returned \texttt{FAIL} in all 5 iterations for matrices characterized by high row density. This discrepancy from the sequential baseline is not due to logical or algorithmic flaws, but to intrinsic non-associativity of single-precision floating-point arithmetic (\texttt{Float32}). In a parallel execution context, non-deterministic accumulation order introduces unavoidable rounding errors. On rows with many non-zero elements, these errors accumulate until generating numerical drift that exceeds the predefined fixed tolerance threshold ($10^{-3}$).


\section{Discussion}TTS analysis highlights a fundamental aspect: the enormous computational speed of GPUs compared to CPUs is often nullified by the time required to read files from disk and transfer data through the PCIe bus. This means optimizing CUDA code is crucial if multiplication is repeated many times (as in iterative algorithms like Conjugate Gradient), where the cost of I/O is amortized over thousands of iterations. Conversely, in "one-shot" computations, the total time depends almost entirely on the latency of preparatory steps and the overhead of host-to-device memory transfers.Matrix structure remains the determining factor for performance. In irregular matrices, the Base CSR kernel fails due to severe \textit{warp divergence}: threads within the same warp process rows of vastly different lengths, forcing faster threads to remain idle until the slowest finishes. This is solved by the CSR-Vector kernel, which assigns an entire warp to a single row. This approach enables memory coalescing, allowing the hardware to group multiple memory requests into a single transaction.Despite this, throughput remains far from the theoretical peak of the NVIDIA A30. This gap is primarily caused by indirect memory accesses ($x[col\_idx[j]]$): while matrix values are fetched contiguously, the elements of the input vector $x$ are accessed following the sparsity pattern, causing cache misses and fragmented data transfers. In this scenario, the cuSPARSE library proves superior as it employs heuristics to dynamically adapt its execution model to the specific matrix topology, effectively balancing load distribution and memory throughput.
\section{Conclusion}

This study confirms that SpMV efficiency on GPUs is tightly constrained by data topology and memory bandwidth limits.

\begin{itemize}

    \item Key Results:The CSR-Vector kernel proved essential for handling irregular matrices, overcoming divergence limitations of scalar CSR. However, infrastructural costs (I/O and PCIe) represent the real time limitation in single executions.

    \item Practical Implications: For real-world applications, the use of cuSPARSE is recommended due to its dynamic adaptability. Custom kernel development must carefully consider the trade-off between numerical precision (\texttt{Float32}) and speed, especially on high-density datasets.

    \item Future Work: Further investigations could explore the use of mixed precision and configurable tolerance parameters to mitigate numerical drift observed during validation.

\end{itemize}


\begin{thebibliography}{00}

\bibitem{bGao2024} J. Gao, B. Liu, W. Ji, and H. Huang, ``A Systematic Literature Survey of Sparse Matrix-Vector Multiplication,'' \textit{arXiv preprint arXiv:2404.06047}, 2024.

\bibitem{bChu2023} G. Chu, Y. He, L. Dong, Z. Ding, D. Chen, H. Bai, x. Wang, and C. Hu, ``Efficient Algorithm Design of Optimizing SpMV on GPU,'' in \textit{Proc. 32nd Int. Symp. High-Perform. Parallel Distrib. Comput. (HPDC '23)}, 2023, pp. 243–256. doi: 10.1145/3588195.3593002.

\bibitem{bBell2009} N. Bell and M. Garland, ``Implementing sparse matrix-vector multiplication on throughput-oriented processors,'' in \textit{Proc. Conf. High Perform. Comput. Netw. Storage Anal. (SC '09)}, 2009, pp. 1–11.

\bibitem{bMerrill2016} D. Merrill and M. Garland, ``Merge-based parallel sparse matrix-vector multiplication,'' in \textit{Proc. Int. Conf. High Perform. Comput. Netw. Storage Anal. (SC '16)}, 2016, pp. 678–689.

\end{thebibliography}

\end{document}